\chapter{Género, sexualidad y representación}
El Renacimiento replanteó los roles de género en diálogo con nuevas formas de sociabilidad. Saslow explora cómo la iconografía de Ganímedes reveló debates sobre deseo masculino y autoridad divina, proponiendo lecturas que oscilan entre la idealización y la condena moral.\cite{Saslow1986} Estas imágenes funcionaron como campos de prueba para los límites de la tolerancia social.

\begin{figure}[H]
    \centering
    % Reemplazar por la imagen recortada correspondiente
    \includegraphics[width=0.7\textwidth]{FIGURES/ganymede_placeholder}
    \caption{Ganímedes ascendiendo al Olimpo ("Saslow, James M. - Ganymede in the Renaissance. Homosexuality in Art and Society (1986, Yale University Press) - libgen.li.pdf", p. XX).}
    \label{fig:ganymede}
\end{figure}

Las prácticas de autorrepresentación descritas por Woods-Marsden muestran que la identidad del artista también estuvo atravesada por categorías de género, ya que el control de la propia imagen suponía negociar expectativas sociales y religiosas.\cite{WoodsMarsden1998} En paralelo, las narrativas sobre enfermedad venérea estudiadas por Arrizabalaga incorporaron discursos moralizantes que reforzaron normas de comportamiento sexual en las ciudades.\cite{Arrizabalaga1997}
